\subsection{Policy Expressibility and Matched Implementations}
\label{sec:policy-mechanism}

The preceding headline speedups compare different policies: they reflect policy
selection together with the controls needed to execute it, not an intrinsic
speedup from running the same algorithm as BPF. \Cref{tab:policy-expressibility}
distinguishes decision logic from its actuator. User space suffices when it
exposes both the required observations and actions. BPF does not create missing
expert semantics, storage transport or hardware isolation.

\begin{table*}[t]
\centering
\small
\begin{tabularx}{\textwidth}{@{}p{0.20\textwidth}XXX@{}}
\toprule
Decision class & User space & Driver modification & Current \sys support \\
\midrule
Expert prediction/prefetch & Feasible with framework routing and transfer APIs & Not needed for application-owned transfers & Host-JIT component ports; framework adapters supply semantics and transfers \\
Priority queue admission & Feasible with XSched's cooperative Level-1 executor & Not needed for this Level-1 path & Host-JIT HPF decisions; original XSched suspends/resumes queues \\
TSG timeslice control & Our C controller uses a scoped custom driver API & Required for that API and policy hooks & Kernel BPF at the authorized control boundary; native GSP actuator \\
Fault-time prefetch & Explicit application prefetch is possible, but is not a fault-handler policy & Needed to expose the fault decision & Bounded current-VA-block selection \\
UVM eviction-list order & Not exposed by ordinary application hints & Needed to expose PMM list decisions & Validated used/unused head/tail requests; not arbitrary victim eviction \\
SM-local task selection & Feasible inside a cooperating fused kernel & Not needed for that kernel's task queue & Device adapter built; actual POD comparison remains pending \\
Disk/CXL placement & Storage frameworks can issue their own I/O & New destination/lifetime/completion interfaces needed for driver-managed tiers & No general storage/CXL transition backend \\
\bottomrule
\end{tabularx}
\caption{Expressibility depends on observations and actuators. Host-JIT ports,
kernel-driver policies and device adapters are different paths. A source mapping
or built adapter is not a completed performance or safety test.}
\label{tab:policy-expressibility}
\end{table*}

\paragraph{Matched-policy methodology}
The four application-policy studies below share the workload within each
experiment, and the matched native/BPF pair shares its executor and actuator.
They use RTX 5090, Linux 6.15.11, NVIDIA 575.57.08 and a fixed 400\,W cap, with
exclusive GPU execution, randomized complete blocks and numerical/engagement
checks. Central values are medians across cells unless stated otherwise;
ratios and percentage changes are paired geometric-mean effects, except that
XSched uses means of paired absolute or relative differences. Intervals are paired
whole-block 95\% bootstrap intervals, not intervals on
pooled requests. These studies use explicit compatibility/framework adapters,
and do not establish transparent application integration or reproduction of
the papers' original hardware and headline speedups.

\paragraph{Driver-path mechanism cost}
A separate RTX 5090 experiment on NVIDIA 610.43.02 holds the no-prefetch decision
constant: a native driver switch versus a kernel-verified BPF callback and
checked region setter. All 15 paired blocks read one word per 64\,KiB region
of an 8\,GiB CPU-initialized allocation, with 3,932,160 exact output checks in
total. An untimed migration trace confirms zero actual prefetch in both arms.
BPF increases CUDA-event kernel time by 3.219\% [2.247\%, 4.202\%]. This is a
measured mechanism cost for CPU-resident, non-first-touch UVM faults, not an
application speedup or a result for the different 575 stack. Native and BPF
cell medians are 364.041 and 374.402\,ms; descriptive fault-span rate is not
payload bandwidth. This experiment illustrates why generality costs must be
reported separately from gains due to a better policy.

\paragraph{MoE-Infinity}
We port MoE-Infinity v3's activation-matrix matching, prefetch ranking and
eviction rule~\cite{xue2025moeinfinity} to a common GPT-OSS-120B frontend. Native
and BPF ports share prediction-set protection and the repaired executor.
The prediction-off baseline retains the upstream prefill extra-expert overload
shortcut, whereas both policy ports enforce the strict pool budget. Its
comparison therefore changes policy and executor behavior, not prediction alone.
Five blocks (15 cells) give 11.8964, 11.2233 and 11.1900 token/s for prediction-off,
native policy and BPF. BPF/native throughput is 0.996540 [0.989239, 1.005508],
which establishes neither superiority nor formal equivalence. BPF loses 6.92\%
throughput versus prediction-off while reducing first-visible-text latency by
22.48\%. This combined policy/executor configuration trades throughput for
earlier output on this workload. Throughput includes the final prefetch drain.
Its host uBPF selectors are not a kernel-UVM policy, and this end-to-end
comparison does not isolate JIT instruction cost.

\paragraph{XSched and GPreempt}
\Cref{fig:matched-scheduling} separates foreground protection from background
progress. For XSched~\cite{shen2025xsched}, queue-entry p99 spans host submission
to first CTA entry, including waiting; background throughput counts completed
kernels. BPF implements highest-priority-first decisions over XSched queue state;
the original Level-1 executor performs suspend/resume. Across ten four-arm
blocks, the three displayed arms have p99 medians of 76.8780, 26.9784 and
27.2502\,s. BPF-minus-XSched paired mean change is $+0.292$\,s
$[-0.109,+0.642]$. The separate driver-timeslice arm uses a different policy:
its lower queueing delay costs 39.5\% of XSched's background throughput
(paired-mean relative change) and is
not a substitute for the matched HPF comparison.

For GPreempt~\cite{fan2025gpreempt}, response p99 spans scheduled arrival through
GPU-synchronized, numerically verified completion. Background goodput counts
only verified completions inside the 60-second window. VGG19 foreground runs at
100 requests/s against ResNet152 at 100, 200 or continuous background supply;
each load has five three-arm blocks. The models use seeded FP32 TVM testing
weights, batch-one CUDA graphs and 200\,$\mu$s preprocessing. Native uses one
prioritized-stream context; C/BPF share a two-context executor on the same custom
driver. BPF executes host-JIT reset/hint/block/release decisions and kernel-BPF
timeslice callbacks. C and BPF share mapped-host flags because
original GDRCopy transport did not work on this GPU. Both improve foreground
p99 over native at all three loads. With continuous background work, C/native
reduces p99 by 9.61\% [6.46\%, 11.91\%] and BPF/native by 10.81\%
[9.91\%, 11.70\%], at approximately 9\% background-goodput cost. BPF/C p99 changes
by $-1.32$\% $[-4.19\%,+0.51\%]$. All 270,000 offered foreground requests completed
correctly within their windows. Neither latency metric measures a hardware
preemption quantum; the two workloads do not provide a cross-system ranking.

\paragraph{FineMoE dynamic prefetch sets}
The FineMoE~\cite{yu2026finemoe} component experiment compares demand-only,
an all-positive prefetch ablation, the native dynamic-set rule and its host-BPF port on
Qwen1.5-MoE-A2.7B-Chat with a shared offloader. Five blocks use eight held-out
public MT-Bench prompts, at most 16 input tokens and exactly 16 generated tokens;
the full gated LMSYS dataset was unavailable. The numerical preflight checks
87,515,136 saved FP32 values exactly; all 2,560 formal tokens match the original
model. Unlike the MoE-Infinity metric, this throughput excludes final drain.
BPF reduces completed prefetch payload evicted before demand use by
60.41\% versus all-positive and increases throughput by 24.64\%. However, BPF
throughput is 12.62\% lower than demand-only. BPF/native throughput is
1.002107 [0.998266, 1.005226]. These are logical completed payload bytes, not
physical PCIe traffic; resident-but-unused copies are not classified as wasted.
The benefit over all-positive comes from the selection rule, not BPF, and
speculation still loses aggregate throughput to demand-only.

\paragraph{Policy origin and boundaries}
These ports implement existing algorithms, not agent-discovered policies.
The original case studies explore compositions such as stride prefetch plus LFU,
region-specific prefetch and phase-sensitive admission; generation by an agent
does not by itself establish algorithmic novelty. In
\Cref{fig:all-kernels-priority}, historical single-round memory results motivate
resource-specific policies, but unverified scheduler engagement prevents a claim
that a correctly running scheduler was ineffective. A same-policy native/BPF
comparison instead measures implementation cost. Host-JIT
numerical parity does not establish kernel-verifier admission or SIMT-verifier
enforcement; the latter's deployment limits are explicit in
\S\ref{sec:verification}.
% TODO: Add complete LMCache local-disk, qualified Expert Buffering analogue,
% and final Hummingbird/POD results. Preparation and incomplete cells do not count.
